\section*{第 \ref{ch:b1:induction} 章——电磁感应及其应用}

\begin{solution}{exo:b1:induction:1}
$\Phi = NBS\cos30^\circ = 100 \times 0.2 \times 10^{-3} \times 0.866 = \qty{1.7e-2}{Wb}$;
$\langle e\rangle = \Delta\Phi/\Delta t = \qty{1.7}{V}$ 在 \qty{10}{ms} 内，\qty{17}{V}，
在 \qty{1}{ms} 内。电流方向使其磁场重建正在消失的磁通量：
其磁场沿原来的 $\vect B$ 方向。
\end{solution}

\begin{solution}{exo:b1:induction:2}
$e = B\ell v = 0.5 \times 0.2 \times 2 = \qty{0.20}{V}$; $i = \qty{2.0}{A}$; $F = B\ell i =
\qty{0.20}{N}$，为制动力，故施加的推力为 \qty{0.20}{N}；$Fv = \qty{0.40}{W}
= Ri^2$.
\end{solution}

\begin{solution}{exo:b1:induction:3}
$L = \mu_0n^2S\ell = 4\pi \times 10^{-7} \times 10^6 \times 4 \times 10^{-4} \times 0.2 =
\qty{0.10}{mH}$; $W = \tfrac12LI^2 = \qty{1.3}{mJ}$; $\tau = L/R = \qty{0.1}{ms}$.
$B = \mu_0nI = \qty{6.3}{mT}$，$B^2/2\mu_0 = \qty{15.7}{J/m^3}$，乘以
体积 $S\ell = \qty{8e-5}{m^3}$，得 \qty{1.3}{mJ}。
\end{solution}

\begin{solution}{exo:b1:induction:4}
$N_2/N_1 = 12/230 = 0.052$；$N_1 = 50/0.052 = 960$ 匝；$i_1 = 2 \times 0.052 =
\qty{0.10}{A}$。恒定电流产生恒定磁通，故不存在 $\dd\varphi/\dd t$，
次级无感应\omterm{def:b1:dc-circuits:voltage}{电压}——而初级相当于短路，会被烧毁。
\end{solution}

\begin{solution}{exo:b1:induction:5}
(a) 沿接近方向的磁通量增加：感应电流产生磁场
与之相反，其自身北极正对磁铁北极——而圆环
被推开。(b) 感应电流的磁场阻碍磁通量增加；
圆环相当于与 $\vect B$ 反向排列的偶极子，被推向较弱
磁场处。(c) 下落时磁通量增加：电流与 (a) 相同，圆环
被制动。(d) 磁铁前方的管段所见磁通量增加，
后方则减少；两段感应电流都阻碍磁铁运动。在
速率达到使\omterm{thm:b1:magnetostatics:laplace}{拉普拉斯力}与重力平衡时，磁铁匀速下落，
其失去的\omterm{def:b1:work-and-energy:conservative}{势能}转化为管中的焦耳热。
\end{solution}

\begin{solution}{exo:b1:induction:6}
$E_{\max} = NBS\omega = 200 \times 2 \times 10^{-3} \times 0.1 \times 314 = \qty{12.6}{V}$，\omterm{def:b1:sinusoidal-impedance:rms}{有效值}
\qty{8.9}{V}; $I_{\max} = \qty{1.26}{A}$, $\Gamma_{\max} = NBSI_{\max} = \qty{0.050}{N.m}$
（当 $\sin\omega t = 1$ 时）；$\langle P\rangle = E_{\max}^2/2R = \qty{7.9}{W}$ $=
\langle\Gamma\omega\rangle = \tfrac12 \times 0.05 \times 314$；当 $e = 0$ 时，电流为
零，力矩随之消失——力矩以 \qty{100}{Hz} 的频率脉动。
\end{solution}

\begin{solution}{exo:b1:induction:7}
$e = Bav$, $i = Bav/R$, $F = B^2a^2v/R = 10\,v$ (N). $m\,\dd v/\dd t = -B^2a^2v/R$:
$v = v_0\eu^{-t/\tau}$，$\tau = mR/B^2a^2 = \qty{0.10}{s}$；距离 $v_0\tau =
\qty{10}{cm}$（若磁场区域足够长）。\omterm{thm:b1:work-and-energy:ke}{动能}
\qty{0.5}{J}，以 $Ri^2$ 的形式在回路中耗散。
\end{solution}

\begin{solution}{exo:b1:induction:8}
螺线管磁场 $\mu_0n_1i_1$ 穿过 $N_2$ 匝线圈的磁通量：$M =
\mu_0n_1N_2S_2 = 4\pi \times 10^{-7} \times 2000 \times 50 \times 2 \times 10^{-4} = \qty{25}{\micro H}$.
电流线性变化时：$e = M\,\dd i_1/\dd t = \qty{2.5}{mV}$；在 \qty{10}{kHz}、\qty{1}{A} 时：$e_{\max} =
M\omega I = 25 \times 10^{-6} \times 6.3 \times 10^4 = \qty{1.6}{V}$。$M$ 是对称的：
同一数值也给出小线圈穿过
螺线管的磁通量——直接计算要困难得多。
\end{solution}

\begin{solution}{exo:b1:induction:9}
$W = B^2V/2\mu_0 = 2.25/(2.51 \times 10^{-6}) = \qty{0.90}{MJ}$; $L = 2W/I^2 =
\qty{7.2}{H}$；$h = W/mg = \qty{92}{m}$；氦气：$9 \times 10^5/2600 = \qty{350}{L}$。
地磁场：$(5 \times 10^{-5})^2/2\mu_0 = \qty{1e-3}{J/m^3}$，房间内共
\qty{50}{mJ}。
\end{solution}

\begin{solution}{exo:b1:induction:10}
(a) $E - B\ell v = Ri$（反电动势与电源相反），$m\,\dd v/\dd t =
B\ell i$: $m\,\dd v/\dd t = B\ell(E - B\ell v)/R$. (b) $v_\infty = E/B\ell$, $\tau =
mR/B^2\ell^2$。(c) 将电路方程乘以 $i$：$Ei = Ri^2 +
B\ell vi$，且 $B\ell vi = Fv$ 为机械功率；\omterm{def:b1:heat-engines:efficiency}{效率} $Fv/Ei =
B\ell v/E = v/v_\infty$——仅在空载时为 $100\%$，此时无
输出。(d) $v_\infty = 2/(0.5 \times 0.2) = \qty{20}{m/s}$，$\tau = 0.05 \times
0.1/0.01 = \qty{0.5}{s}$. (e) $F = B\ell I = 1 \times 0.1 \times 10^6 = \qty{1e5}{N}$
（取 $\ell = \qty{10}{cm}$），$a = \qty{1e7}{m/s^2}$，$v = \sqrt{2a \times 5} =
\qty{1e4}{m/s}$——即每秒十千米，这既是
轨道炮的希望，也是其难题（导轨会被侵蚀）。
\end{solution}

\begin{solution}{exo:b1:induction:11}
(a) 在半径 $r$ 处，金属以 $v = \omega r$ 运动；$\vect v\wedge\vect B$ 沿
径向，大小为 $\omega rB$：$e = \int_0^R\omega Br\,\dd r = \tfrac12B\omega R^2$。
(b) $\omega = 314$：$e = 0.5 \times 1 \times 314 \times 0.01 = \qty{1.6}{V}$。(c) 只有
“一匝”，故\omterm{def:b1:dc-circuits:voltage}{电压}至多为\omterm{thm:b1:potential-capacitors:potential}{伏特}量级；接入 \qty{1}{m\ohm}，电流为 \qty{1.6}{kA}，力矩
$= P/\omega = ei/\omega = 2500/314 = \qty{8}{N.m}$。(d) 回路为电刷–半径--
轮缘–电刷，扫过的面积为 $\tfrac12R^2\omega\,\dd t$（每 $\dd t$ 时间内），故 $\dd\Phi/\dd t =
\tfrac12BR^2\omega$——即运动且形变的回路所穿过的磁通量。
\end{solution}

\begin{solution}{exo:b1:induction:12}
(a) $e_2 = -M\,\dd i_1/\dd t$，峰值 $M\omega I_1 = 10^{-6} \times 1.57 \times 10^5 \times 30 =
\qty{4.7}{V}$。(b) $I_2 = 4.7/0.01 = \qty{470}{A}$ 峰值；$\langle P\rangle = R_pI_2^2/2
= \qty{1.1}{kW}$。(c) 电动势 $\propto\omega$：高频时
即使 $M$ 较小也能获得有用的电动势；玻璃为绝缘体，其中
没有电流流过——热量产生于金属内部。(d) 锅中电流在
线圈中感应出 $M\omega I_2 = 10^{-6} \times 1.57 \times 10^5 \times 470 = \qty{74}{V}$ 峰值，
与 $i_1$ 同相（锅中电流落后 $e_2$ 的\omterm{def:b1:wave-propagation:signal}{相位}为零，因其
为阻性，而 $e_2$ 落后 $i_1$ \qty{90}{\degree}\ldots\ 故
反映到线圈中的电动势为 $-M\,\dd i_2/\dd t$，与 $i_1$ 的
\emph{驱动}\omterm{def:b1:dc-circuits:voltage}{电压}反相）：发电机额外提供 $\tfrac12 \times 74 \times 30 =
\qty{1.1}{kW}$，恰好等于锅所耗散的功率。
\end{solution}

\begin{solution}{pb:b1:induction:1}
\textbf{1.} $\Phi = NBS\cos\omega t$, $e = NBS\omega\sin\omega t$; $NBS = 300 \times
0.3 \times 3 \times 10^{-4} = \qty{0.027}{Wb}$; $\omega = v/r = 5/0.01 = \qty{500}{rad/s}$:
\omterm{def:b1:wave-propagation:signal}{振幅} \qty{13.5}{V}，频率 $\omega/2\pi = \qty{80}{Hz}$。

\textbf{2.} $I_{\text{rms}} = 13.5/\sqrt2/16 = \qty{0.60}{A}$；灯泡 $RI^2 = \qty{4.3}{W}$
在 \qty{7.2}{V} 下，一只 \qty{6}{V}、\qty{3}{W} 的灯泡过载，
不能持久。

\textbf{3.} 机械功率 $=$ \omterm{thm:b1:dc-circuits:kirchhoff}{电功率}：$(R + r_c)I_{\text{rms}}^2 = 16 \times 0.36 =
\qty{5.7}{W}$；$F = P/v = 5.7/5 = \qty{1.1}{N}$——约为滚动摩擦力的一半；
骑车人能感觉到。

\textbf{4.} 感应电流所受的\omterm{thm:b1:magnetostatics:laplace}{拉普拉斯力}偶阻碍转动
（楞次）；灯泡的能量来自骑车人的腿，共 \qty{5.7}{W}。

\textbf{5.} $\underline Z = R + r_c + jL\omega$；$I = NBS\omega/\sqrt{(R + r_c)^2 + L^2\omega^2}$；
当 $L\omega \gg R + r_c$ 时，$I \to NBS/L = 0.027/0.02 = \qty{1.35}{A}$。$70\%$：$L\omega
= 0.7(R + r_c)/\sqrt{1 - 0.49} = 15.7$，$\omega = \qty{784}{rad/s}$，$v = \qty{7.8}{m/s}
= \qty{28}{km/h}$。

\textbf{6.} 不计 $L$ 时：$e_{\max} = 3 \times 13.5 = \qty{40.5}{V}$，$I_{\text{rms}} =
\qty{1.8}{A}$，灯泡中功率 \qty{39}{W}——灯泡烧毁。计入 $L$ 时：$I = 40.5/\sqrt{256 +
900} = \qty{1.2}{A}$（峰值），功率 \qty{8.5}{W}：灯泡发热，但小了十倍。

\textbf{7.} \qty{9}{km/h}：$e_{\max} = 6.75$，$|Z| = \sqrt{256 + 25} = 16.8$，$I = 0.40$，
$P = \tfrac12 \times 0.16 \times 12 = \qty{1.0}{W}$：灯光昏暗。\qty{36}{km/h}：$e_{\max} = 27$，
$|Z| = 25.6$，$I = 1.05$，$P = \qty{6.6}{W}$。速率为原来的四倍，
电流为 $2.6$ 倍：电感起\omterm{def:b1:optical-instruments:eye}{调节}作用。

\textbf{8.} 只有\omterm{def:b1:non-inertial-frames:frames}{相对运动}才有影响：穿过固定
线圈的磁通量变化完全相同，电动势也相同。不需要滑动触点（电刷）
来把转动线圈中的电流引出。

\textbf{9.} 无运动、无磁通量变化、无灯光——因此\omterm{def:b1:potential-capacitors:capacitance}{电容器}或
小电池（“驻车灯”）在骑行时充电。

\textbf{10.} 每段导线都垂直于径向的 $\vect B$：
$F = B\ell i$，沿轴向。以速度 $v$ 运动时，每段导线都感受到 $\vect v\wedge\vect B$
沿导线方向：$e = -B\ell v$，与产生
运动的电流方向相反（\omterm{thm:b1:dc-circuits:kirchhoff}{受电器约定}：\omterm{def:b1:dc-circuits:voltage}{电压}为 $B\ell v$）。

\textbf{11.} $u = Ri + L\,\dd i/\dd t + B\ell v$；$m\,\dd v/\dd t = -kx - hv + B\ell i$。

\textbf{12.} $\underline i = (\underline u - B\ell\underline v)/(R + jL\omega)$；$\underline v\,(jm\omega + h +
k/j\omega) = B\ell\,\underline i$；代入并解出 $\underline v$。

\textbf{13.} 由 $i \approx (u - B\ell v)/R$ 可得 $m\dot v + (h + B^2\ell^2/R)v + kx =
B\ell u/R$：附加阻尼 $25/8 = \qty{3.1}{N.s/m}$。$f_0 = \tfrac1{2\pi}\sqrt{k/m} =
\qty{50}{Hz}$；$Q = \sqrt{km}/h_{\text{总}}$：开路时 $3.2$，接
入放大器后为 $0.77$。

\textbf{14.} $F = \qty{5.0}{N}$，$a = F/m = \qty{500}{m/s^2}$，$x = a/\omega^2 =
\qty{13}{\micro m}$，$v = a/\omega = \qty{0.080}{m/s}$；$B\ell v = \qty{0.40}{V}$，与
$Ri = \qty{8}{V}$ 相比为 $5\%$。

\textbf{15.} 若电流给定，则 $v = F/h = \qty{5}{m/s}$（
\omterm{def:b1:wave-propagation:signal}{振幅} \qty{16}{mm}——对 \qty{1}{A} 的电流而言位移很大）；$B\ell v =
\qty{25}{V}$，为 $Ri$ 的三倍：放大器须提供 \qty{33}{V} 以推动
\qty{1}{A} 的电流（共振时）——即\omterm{def:b1:sinusoidal-impedance:impedance}{阻抗}峰值。\omterm{def:b1:dc-circuits:sources}{电压源}型放大器
则会改由电磁阻尼稳住纸盆。

\textbf{16.} $\tfrac12RI^2 = \qty{4.0}{W}$；$\tfrac12hv^2 = \tfrac12 \times 0.0064 =
\qty{3.2}{mW}$：\omterm{def:b1:heat-engines:efficiency}{效率} $0.08\%$——扬声器其实是
会轻声细语的加热器。

\textbf{17.} 共振以上，$a = F/m = B\ell i/m$ 在给定电流下与频率无关，
故压强 $\propto a$ 为平直的：
质量控制频段即为工作频段。

\textbf{18.} 径向磁场：力沿轴向，且
线圈在气隙中任一位置处都相同。高音单元须保持质量控制
到 \qty{20}{kHz}，位移极小：线圈与纸盆须轻；同时
小纸盆会向各方向辐射高频。

\textbf{19.} 低频时：$v$ 很小，$|Z| \approx R = \qty{8}{\ohm}$；在
\qty{50}{Hz}：$u = Ri + B\ell v$，其中 $v = B\ell i/h$：$Z = R + B^2\ell^2/h = 8 + 25 =
\qty{33}{\ohm}$；在 \qty{10}{kHz}：$|R + jL\omega| = |8 + 31j| = \qty{32}{\ohm}$。
共振处出现峰值，平坦的 \qty{8}{\ohm} 基线，随 $L\omega$ 增大而上升。

\textbf{20.} $e = B\ell v = 5 \times 10^{-3} = \qty{5}{mV}$。

\textbf{21.} $i = 5 \times 10^{-3}/608 = \qty{8.2}{\micro A}$；$F = B\ell i = \qty{41}{\micro N}$，
与 $v$ 方向相反：$F = [B^2\ell^2/(R + R_L)]\,v$，相当于阻尼 \qty{0.041}{N.s/m}；
功率 $Fv = \qty{4e-8}{W}$ 取自声音，最终
消耗在两个\omterm{def:b1:dc-circuits:resistor}{电阻器}中。

\textbf{22.} 线圈骨架：$B^2\ell_f^2/R_f = 1 \times 0.01/10^{-3} = \qty{10}{N.s/m}$，为
$h$ 的十倍：会使纸盆发闷，并把功率白白变为热量；开一条缝
即可切断闭合的匝。

\textbf{23.} 开路时：仅 $h$ 起阻尼作用，$Q = 3.2$，会持续振荡。短路时：
反电动势驱动电流流过 $R$，附加阻尼 $B^2\ell^2/R$，
$Q = 0.77$：立即停止振荡。（可自己试一试。）

\textbf{24.} 一个周期内，储能回到原值：
$\langle ui\rangle = \langle Ri^2\rangle + \langle hv^2\rangle$——电输入 $=$
线圈中的焦耳热 $+$ 机械损耗（辐射的声音与
悬挂摩擦力）；耦合项 $B\ell iv$ 在两式中都出现，
且符号相反，相互抵消。

\textbf{25.} 发电机：转动线圈上的\omterm{def:b1:potential-capacitors:capacitance}{法拉}第电磁感应——速率 \qty{18}{km/h} 时
\omterm{def:b1:dc-circuits:voltage}{电压}为 \qty{13.5}{V}，电流受自身电感限制，至多 \qty{1.35}{A}。
扬声器：\omterm{thm:b1:magnetostatics:laplace}{拉普拉斯力}与反电动势——提供 \qty{3.1}{N.s/m} 的电
磁阻尼，\omterm{def:b1:heat-engines:efficiency}{效率} $0.08\%$。麦克风：同一线圈反向使用——
每毫米每秒产生 \qty{5}{mV}。
\end{solution}
